\section{Figure Plan}
\label{sec:figure-plan}

% Use this section as a working plan for the paper's figures. It can be removed,
% moved to an appendix, or converted into actual figure environments once the
% venue and final paper structure are clear.
%
% Suggested notes for each figure:
% - What question does this figure answer?
% - What data, image, schematic, or result will it show?
% - What is the intended takeaway?
% - What still needs to be generated, exported, or redesigned?

\begin{itemize}
\item \textbf{Figure 1:} Overview of the proposed sensor, the two sensing
modalities, and the principal experimental evaluations. Dhruv

\item \textbf{Figure 2:} Sensor architecture and design, including the
material structure, electrode arrangement, fabrication process, and
electrical interfacing. Aditya 

\item \textbf{Figure 3:} Experimental setup and data-collection procedure,
including contact-position control, force application, and signal
acquisition. Aditya

\item \textbf{Figure 4:} Fundamental sensor characterisation, including
signal stability, repeatability, sensitivity, hysteresis, drift, and
operating force range. Dhruv

\item \textbf{Figure 5:} Signal-processing pipelines for the direct
resistive-grid and EIT sensing modalities, including preprocessing,
feature extraction, EIT reconstruction, and localisation. Aditya

\item \textbf{Figure 6:} Spatial sensing performance using direct
resistive-grid measurements, including single-contact localisation,
spatial resolution, force dependence, and multi-contact response. Aditya

\item \textbf{Figure 7:} Spatial sensing performance using EIT
measurements, including single-contact localisation, spatial resolution,
force dependence, and multi-contact response. Aditya and Dhruv

\item \textbf{Figure 8:} Direct comparison of the two sensing modalities,
including localisation accuracy, spatial separability, measurement
superposition, multi-touch performance, and modality-specific failure
modes. Dhruv

\item \textbf{Figure 9:} Optional extension investigating distributed
contact-map estimation from the single-contact localisation dataset.

\item \textbf{Figure 10:} Optional demonstration of multi-contact or simple
contact-shape reconstruction using the estimated virtual-taxel
representation.


\end{itemize}
